\subsection{Adjoints and Fundamental Subspaces}


% **Role:** linear algebra with metric and angle structure.
%
% Adjoints, conjugate transpose, invariant subspaces, kernels, and ranges.

Adjoints express dual action in inner-product language. They become conjugate transposes in orthonormal coordinates and exchange kernels, ranges, and orthogonal complements.

Throughout this subsection, $\mathbb{F}$ is $\mathbb{R}$ or $\mathbb{C}$.

\begin{definition} \label{def:adjoint-linear-map} \hlt{Adjoint of a Linear Map}\\
Let $\mathcal{V}$ and $\mathcal{W}$ be inner-product spaces and let $T\in\LL(\mathcal{V},\mathcal{W})$.\\
A function $S:\mathcal{W}\to\mathcal{V}$ is an \hlt{adjoint} of $T$ if $(\forall\,\mathbf{v}\in\mathcal{V})(\forall\,\mathbf{w}\in\mathcal{W})\,\innerprod[\mathcal{W}]{T\mathbf{v},\mathbf{w}}=\innerprod[\mathcal{V}]{\mathbf{v},S\mathbf{w}}$.
\end{definition}

\begin{perspective} \label{per:dual-map-versus-inner-product-adjoint} \hlt{Dual Map and Inner-Product Adjoint}\\
For $T:\mathcal{V}\to\mathcal{W}$, the pullback $\varphi\mapsto\varphi\circ T$ is the \hlt{dual map} $T':\mathcal{W}'\to\mathcal{V}'$, also called the \hlt{operator adjoint} and written $T^\times$.\\
This algebraic map is distinct from the inner-product adjoint $T^*$, whose construction requires inner products.
\end{perspective}

\begin{lemma} \label{lem:adjoint-unique-if-exists} \hlt{Uniqueness of Adjoints}\\
Let $\mathcal{V}$ and $\mathcal{W}$ be inner-product spaces and let $T\in\LL(\mathcal{V},\mathcal{W})$. Then $T$ has at most one adjoint.
\begin{proof}
Suppose $S_1,S_2:\mathcal{W}\to\mathcal{V}$ are adjoints of $T$, $\mathbf{w}\in\mathcal{W}$. Then $(\forall\,\mathbf{v}\in\mathcal{V})\,\innerprod[\mathcal{V}]{\mathbf{v},S_1\mathbf{w}}=\innerprod[\mathcal{V}]{\mathbf{v},S_2\mathbf{w}}$ by the two adjoint identities.\\
Hence $S_1\mathbf{w}=S_2\mathbf{w}$ by Lemma \ref{lem:inner-products-separate-vectors}(i). Since $\mathbf{w}$ was arbitrary, $S_1=S_2$.\\
\hlt{Structural Commentary}: Pointwise comparison reduces uniqueness to separation by inner products.\\
\hlt{Key Insight}: Distinct vectors cannot agree against every inner-product test.
\end{proof}
\end{lemma}

\begin{theorem} \label{thm:finite-dimensional-adjoint-exists} \hlt{Existence of the Adjoint}\\
Let $\mathcal{V}$ be a finite-dimensional inner-product space, let $\mathcal{W}$ be an inner-product space, and let $T\in\LL(\mathcal{V},\mathcal{W})$. Then:
\begin{enumerate}[label=(\roman*)]
\setlength{\itemsep}{0pt}
\item $T$ has an adjoint $T^*:\mathcal{W}\to\mathcal{V}$;
\item the adjoint $T^*$ is unique;
\item $T^*\in\LL(\mathcal{W},\mathcal{V})$.
\end{enumerate}
\begin{proof}
\begin{enumerate}[label=(\roman*)]
\setlength{\itemsep}{0pt}
\item Fix $\mathbf{w}\in\mathcal{W}$. The map $\varphi_{\mathbf{w}}\defeq\innerprod[\mathcal{W}]{T(\cdot),\mathbf{w}}$ belongs to $\mathcal{V}'$ by linearity of $T$ and the first inner-product slot.\\
The unique $T^*\mathbf{w}\in\mathcal{V}$ satisfying $(\forall\,\mathbf{v}\in\mathcal{V})\,\innerprod[\mathcal{W}]{T\mathbf{v},\mathbf{w}}=\innerprod[\mathcal{V}]{\mathbf{v},T^*\mathbf{w}}$ exists by Theorem \ref{thm:riesz-representation-finite-dimensional}, and resulting map $T^*$ is adjoint.
\item The adjoint $T^*$ is unique by Lemma \ref{lem:adjoint-unique-if-exists}.
\item Fix $a\in\mathbb{F}$, $\mathbf{w}_1,\mathbf{w}_2\in\mathcal{W}$. Then $\innerprod[\mathcal{V}]{\mathbf{v},T^*(a\mathbf{w}_1+\mathbf{w}_2)}=\innerprod[\mathcal{W}]{T\mathbf{v},a\mathbf{w}_1+\mathbf{w}_2}=\overline{a}\innerprod[\mathcal{W}]{T\mathbf{v},\mathbf{w}_1}+\innerprod[\mathcal{W}]{T\mathbf{v},\mathbf{w}_2}=\innerprod[\mathcal{V}]{\mathbf{v},aT^*\mathbf{w}_1+T^*\mathbf{w}_2}$ for $\mathbf{v}\in\mathcal{V}$ by conjugate linearity in the second slot and the adjoint identity.\\
Hence $T^*(a\mathbf{w}_1+\mathbf{w}_2)=aT^*\mathbf{w}_1+T^*\mathbf{w}_2$ by Lemma \ref{lem:inner-products-separate-vectors}(i), so $T^*$ is linear.
\end{enumerate}
\hlt{Structural Commentary}: Riesz representation constructs codomain-indexed vectors; separation verifies that their assignment is linear.\\
\hlt{Key Insight}: Riesz vectors assemble pointwise into the adjoint.
\end{proof}
\end{theorem}

Figure \ref{fig:adjoint-inner-product-bridge} records the inner-product identity that transfers $T$ from the first slot to $T^*$ in the second slot.
\begin{figure}[H]
  \centering
  \begin{tikzpicture}[
  >=Stealth,
  every node/.style={font=\small},
  space/.style={draw=gray!45, rounded corners=3pt, minimum width=3.6cm, minimum height=2.4cm},
  vec/.style={->, blue!75!black, line width=0.85pt},
  adj/.style={->, teal!75!black, line width=0.85pt},
  map/.style={->, gray!70, line width=0.75pt}
]
  \node[space] (V) at (0,0) {};
  \node[anchor=north west] at (-1.6,1.05) {$\mathcal{V}$};
  \node[space] (W) at (5.2,0) {};
  \node[anchor=north west] at (3.6,1.05) {$\mathcal{W}$};

  \coordinate (Ov) at (-0.95,-0.65);
  \coordinate (v) at (0.55,0.55);
  \coordinate (Tsw) at (0.85,-0.25);
  \draw[vec] (Ov) -- (v) node[midway, above left] {$\mathbf{v}$};
  \draw[adj] (Ov) -- (Tsw) node[midway, below] {$T^*\mathbf{w}$};

  \coordinate (Ow) at (4.25,-0.65);
  \coordinate (Tv) at (5.85,0.55);
  \coordinate (w) at (5.95,-0.25);
  \draw[vec] (Ow) -- (Tv) node[midway, above left] {$T\mathbf{v}$};
  \draw[adj] (Ow) -- (w) node[midway, below] {$\mathbf{w}$};

  \draw[map] (1.9,0.45) -- (3.25,0.45) node[midway, above] {$T$};
  \draw[map] (3.25,-0.45) -- (1.9,-0.45) node[midway, below] {$T^*$};

  \node[draw=gray!45, rounded corners=3pt, align=center, inner sep=5pt] at (2.6,-1.85)
    {$\innerprod[\mathcal{W}]{T\mathbf{v},\mathbf{w}}
      =\innerprod[\mathcal{V}]{\mathbf{v},T^*\mathbf{w}}$};
\end{tikzpicture}

  \caption{Adjoint relation: $\innerprod[\mathcal{W}]{T\mathbf{v},\mathbf{w}}=\innerprod[\mathcal{V}]{\mathbf{v},T^*\mathbf{w}}$.}
  \label{fig:adjoint-inner-product-bridge}
\end{figure}

\begin{perspective} \label{per:hilbert-adjoint-riesz-dual-map} \hlt{Inner-Product Adjoints as Riesz-Transported Dual Maps}\\
Let $\mathcal{V}$ and $\mathcal{W}$ be finite-dimensional inner-product spaces, $R_{\mathcal{V}}:\mathcal{V}'\to\mathcal{V}$ and $R_{\mathcal{W}}:\mathcal{W}'\to\mathcal{W}$ be Riesz maps from Theorem \ref{thm:riesz-map-conjugate-isomorphism}.\\
For $T\in\LL(\mathcal{V},\mathcal{W})$, the dual map $T':\mathcal{W}'\to\mathcal{V}'$ and the inner-product adjoint $T^*$ satisfy $T'=(R_{\mathcal{V}})^{-1}\circ T^*\circ R_{\mathcal{W}}$.\\
Indeed, if $\varphi\in\mathcal{W}'$ and $\mathbf{v}\in\mathcal{V}$, then $(T'\varphi)(\mathbf{v})=\varphi(T\mathbf{v})=\innerprod[\mathcal{W}]{T\mathbf{v},R_{\mathcal{W}}\varphi}=\innerprod[\mathcal{V}]{\mathbf{v},T^*R_{\mathcal{W}}\varphi}$.\\
Thus $T'$ pulls back functionals by moving their Riesz vectors with $T^*$.\\
\hlt{Structural Commentary}: Riesz representation conjugates the algebraic dual map into an inner-product operator.\\
\hlt{Key Insight}: Dual maps and adjoints encode the same reversed action.
\end{perspective}

Figure \ref{fig:riesz-dual-hilbert-adjoint-square} displays the commuting square expressed by $R_{\mathcal{V}}T'=T^*R_{\mathcal{W}}$.
\begin{figure}[H]
  \centering
  \begin{tikzpicture}[
  x=0.88cm,
  y=0.88cm,
  >=Stealth,
  every node/.style={font=\small},
  space/.style={draw=gray!55, fill=gray!5, rounded corners=3pt,
    minimum width=1.38cm, minimum height=0.60cm, align=center},
  arrow/.style={->, line width=0.78pt, gray!70!black},
  riesz/.style={->, line width=0.72pt, gray!60},
  hint/.style={font=\scriptsize, text=gray!70!black, inner sep=1pt}
]
  \node[space] (Wp) at (0,1.35) {$\mathcal{W}'$};
  \node[space] (Vp) at (4.15,1.35) {$\mathcal{V}'$};
  \node[space] (W) at (0,0) {$\mathcal{W}$};
  \node[space] (V) at (4.15,0) {$\mathcal{V}$};

  \draw[arrow] (Wp) -- node[midway, above, hint] {$T'$} (Vp);
  \draw[arrow] (W) -- node[midway, below, hint] {$T^*$} (V);
  \draw[riesz] (Wp) -- node[midway, left, hint] {$R_{\mathcal{W}}$} (W);
  \draw[riesz] (Vp) -- node[midway, right, hint] {$R_{\mathcal{V}}$} (V);
\end{tikzpicture}

  \caption{Riesz transports the dual map to the adjoint: $R_{\mathcal{V}}T'=T^*R_{\mathcal{W}}$.}
  \label{fig:riesz-dual-hilbert-adjoint-square}
\end{figure}

\begin{definition} \label{def:conjugate-transpose} \hlt{Conjugate Transpose}\\
For $A=[a_{ij}]\in\operatorname{Mat}(m,n,\mathbb{F})$, the \hlt{conjugate transpose} of $A$ is $A^*=[\overline{a_{ji}}]\in\operatorname{Mat}(n,m,\mathbb{F})$.\\
Over $\mathbb{R}$, the conjugate transpose is the ordinary transpose.
\end{definition}

\begin{theorem} \label{thm:matrix-of-adjoint-is-conjugate-transpose} \hlt{Matrix of the Adjoint}\\
Let $\mathcal{V}$ and $\mathcal{W}$ be finite-dimensional inner-product spaces.\\
Let $\beta=(\mathbf{e}_1,\ldots,\mathbf{e}_n)$ and $\gamma=(\mathbf{f}_1,\ldots,\mathbf{f}_m)$ be orthonormal bases of $\mathcal{V}$ and $\mathcal{W}$, respectively. If $T\in\LL(\mathcal{V},\mathcal{W})$, then $[T^*]_{\beta\leftarrow\gamma}=([T]_{\gamma\leftarrow\beta})^*$.
\begin{proof}
Set $A\defeq[T]_{\gamma\leftarrow\beta}$ and $B\defeq[T^*]_{\beta\leftarrow\gamma}$.\\
For $i\in\{1,\ldots,n\}$ and $j\in\{1,\ldots,m\}$, $A_{j,i}=\innerprod[\mathcal{W}]{T\mathbf{e}_i,\mathbf{f}_j}$ and $B_{i,j}=\innerprod[\mathcal{V}]{T^*\mathbf{f}_j,\mathbf{e}_i}$ by Theorem \ref{thm:matrix-entries-from-orthonormal-codomain}.\\
Hence $B_{i,j}=\overline{\innerprod[\mathcal{V}]{\mathbf{e}_i,T^*\mathbf{f}_j}}=\overline{\innerprod[\mathcal{W}]{T\mathbf{e}_i,\mathbf{f}_j}}=\overline{A_{j,i}}$ by conjugate symmetry and the adjoint identity, so $B=A^*$.\\
\hlt{Structural Commentary}: Orthonormal coordinates convert the adjoint identity into entrywise conjugate transposition.\\
\hlt{Key Insight}: Moving an operator across an inner product becomes conjugate transposition in orthonormal bases.
\end{proof}
\end{theorem}

\begin{proposition} \label{prop:basic-adjoint-identities} \hlt{Basic Adjoint Identities}\\
Let $\mathcal{U}$, $\mathcal{V}$, and $\mathcal{W}$ be finite-dimensional inner-product spaces. Then:
\begin{enumerate}[label=(\roman*)]
\setlength{\itemsep}{0pt}
\item $(aS+T)^*=\overline{a}S^*+T^*$ for $S,T\in\LL(\mathcal{V},\mathcal{W})$ and $a\in\mathbb{F}$;
\item $(SR)^*=R^*S^*$ for $R\in\LL(\mathcal{U},\mathcal{V})$ and $S\in\LL(\mathcal{V},\mathcal{W})$;
\item $(T^*)^*=T$ for $T\in\LL(\mathcal{V},\mathcal{W})$;
\item $I_{\mathcal{V}}^*=I_{\mathcal{V}}$.
\end{enumerate}
\begin{proof}
Each candidate is identified by Lemma \ref{lem:adjoint-unique-if-exists} after verification of the adjoint identity.
\begin{enumerate}[label=(\roman*)]
\setlength{\itemsep}{0pt}
\item For $\mathbf{v}\in\mathcal{V}$ and $\mathbf{w}\in\mathcal{W}$, $\innerprod[\mathcal{W}]{(aS+T)\mathbf{v},\mathbf{w}}=a\innerprod[\mathcal{W}]{S\mathbf{v},\mathbf{w}}+\innerprod[\mathcal{W}]{T\mathbf{v},\mathbf{w}}=\innerprod[\mathcal{V}]{\mathbf{v},(\overline{a}S^*+T^*)\mathbf{w}}$. Hence $(aS+T)^*=\overline{a}S^*+T^*$.
\item For $\mathbf{u}\in\mathcal{U}$ and $\mathbf{w}\in\mathcal{W}$, $\innerprod[\mathcal{W}]{SR\mathbf{u},\mathbf{w}}=\innerprod[\mathcal{V}]{R\mathbf{u},S^*\mathbf{w}}=\innerprod[\mathcal{U}]{\mathbf{u},R^*S^*\mathbf{w}}$. Hence $(SR)^*=R^*S^*$.
\item Conjugate symmetry transforms $(\forall\,\mathbf{v}\in\mathcal{V})(\forall\,\mathbf{w}\in\mathcal{W})\,\innerprod[\mathcal{W}]{T\mathbf{v},\mathbf{w}}=\innerprod[\mathcal{V}]{\mathbf{v},T^*\mathbf{w}}$ into $(\forall\,\mathbf{w}\in\mathcal{W})(\forall\,\mathbf{v}\in\mathcal{V})\,\innerprod[\mathcal{V}]{T^*\mathbf{w},\mathbf{v}}=\innerprod[\mathcal{W}]{\mathbf{w},T\mathbf{v}}$. Thus $T$ is the adjoint of $T^*$, so $(T^*)^*=T$.
\item For $\mathbf{v},\mathbf{w}\in\mathcal{V}$, $\innerprod{I_{\mathcal{V}}\mathbf{v},\mathbf{w}}=\innerprod{\mathbf{v},I_{\mathcal{V}}\mathbf{w}}$. Thus $I_{\mathcal{V}}^*=I_{\mathcal{V}}$.
\end{enumerate}
\hlt{Structural Commentary}: Each algebraic operation is transported through the two inner-product slots and then settled by uniqueness.\\
\hlt{Key Insight}: Adjoint formation conjugates scalars, reverses products, and is involutive.
\end{proof}
\end{proposition}

\begin{proposition} \label{prop:adjoints-inverses-and-polynomials} \hlt{Adjoints of Inverses and Polynomials}\\
Let $T\in\LL(\mathcal{V})$ on a finite-dimensional inner-product space. For $p(z)=\sum\limits_{j=0}^{m}a_jz^j\in\mathbb{F}[z]$, set $p^{\#}(z)\defeq\sum\limits_{j=0}^{m}\overline{a_j}z^j$.
\begin{enumerate}[label=(\roman*)]
\setlength{\itemsep}{0pt}
\item If $T$ is invertible, then $(T^*)^{-1}=(T^{-1})^*$.
\item $p(T)^*=p^{\#}(T^*)$.
\item $p\in\mathbb{R}[z]\Longrightarrow p(T)^*=p(T^*)$.
\end{enumerate}
\begin{proof}
\begin{enumerate}[label=(\roman*)]
\setlength{\itemsep}{0pt}
\item Identities $(T^{-1})^*T^*=(TT^{-1})^*=I_{\mathcal{V}}$ and $T^*(T^{-1})^*=(T^{-1}T)^*=I_{\mathcal{V}}$ hold by Proposition \ref{prop:basic-adjoint-identities}(ii),(iv); $(T^*)^{-1}=(T^{-1})^*$.
\item The base identity $(T^0)^*=I_{\mathcal{V}}=(T^*)^0$ holds by Proposition \ref{prop:basic-adjoint-identities}(iv).\\
If $(T^j)^*=(T^*)^j$, then $(T^{j+1})^*=(T^jT)^*=T^*(T^j)^*=(T^*)^{j+1}$ by Proposition \ref{prop:basic-adjoint-identities}(ii).\\
Thus $(\forall\,j\in\{0,\ldots,m\})\,(T^j)^*=(T^*)^j$ by induction.\\
The identity $p(T)^*=(\sum\limits_{j=0}^{m}a_jT^j)^*=\sum\limits_{j=0}^{m}\overline{a_j}(T^*)^j=p^{\#}(T^*)$ holds by repeated application of Proposition \ref{prop:basic-adjoint-identities}(i).
\item If $p\in\mathbb{R}[z]$, then $p^{\#}=p$, so $p(T)^*=p(T^*)$ by (ii).
\end{enumerate}
\hlt{Structural Commentary}: The product law controls inverses and powers, while conjugate additivity controls polynomial coefficients.\\
\hlt{Key Insight}: Adjoint formation applies coefficientwise conjugation to polynomial expressions in an operator.
\end{proof}
\end{proposition}

\begin{proposition} \label{prop:adjoint-kernel-range-orthogonality} \hlt{Adjoint Kernel-Range Orthogonality}\\
Let $\mathcal{V}$ and $\mathcal{W}$ be finite-dimensional inner-product spaces and let $T\in\LL(\mathcal{V},\mathcal{W})$. Then:
\begin{enumerate}[label=(\roman*)]
\setlength{\itemsep}{0pt}
\item $\ker(T^*)=(\operatorname{range}(T))^{\perp}$;
\item $\ker(T)=(\operatorname{range}(T^*))^{\perp}$;
\item $\operatorname{range}(T^*)=(\ker(T))^{\perp}$;
\item $\operatorname{range}(T)=(\ker(T^*))^{\perp}$.
\end{enumerate}
Equivalently, an adjoint annihilates exactly the vectors orthogonal to the range of the original map.
\begin{proof}
\begin{enumerate}[label=(\roman*)]
\setlength{\itemsep}{0pt}
\item For $\mathbf{w}\in\mathcal{W}$, $\mathbf{w}\in\ker(T^*)\Longleftrightarrow(\forall\,\mathbf{v}\in\mathcal{V})\,\innerprod[\mathcal{V}]{\mathbf{v},T^*\mathbf{w}}=0\Longleftrightarrow(\forall\,\mathbf{v}\in\mathcal{V})\,\innerprod[\mathcal{W}]{\mathbf{w},T\mathbf{v}}=0\Longleftrightarrow\mathbf{w}\in(\operatorname{range}(T))^{\perp}$ by Lemma \ref{lem:inner-products-separate-vectors}(ii), the adjoint identity, and conjugate symmetry.
\item The identity $\ker(T)=\ker((T^*)^*)=(\operatorname{range}(T^*))^{\perp}$ holds by (i) for $T^*$ and Proposition \ref{prop:basic-adjoint-identities}(iii).
\item Taking orthogonal complements in (ii), $(\ker(T))^{\perp}=((\operatorname{range}(T^*))^{\perp})^{\perp}=\operatorname{range}(T^*)$ by Corollary \ref{cor:double-orthogonal-complement}.
\item Taking orthogonal complements in (i), $(\ker(T^*))^{\perp}=((\operatorname{range}(T))^{\perp})^{\perp}=\operatorname{range}(T)$ by Corollary \ref{cor:double-orthogonal-complement}.
\end{enumerate}
\hlt{Structural Commentary}: The kernel formula, involution, and double complementation establish all four identities.\\
\hlt{Key Insight}: Kernels and ranges are exchanged by taking an adjoint and an orthogonal complement.
\end{proof}
\end{proposition}

\begin{proposition} \label{prop:adjoint-product-kernels-ranges-and-projections} \hlt{Adjoint Products and Algebraic Projections}\\
Let $\mathcal{V}$ and $\mathcal{W}$ be finite-dimensional inner-product spaces and let $T\in\LL(\mathcal{V},\mathcal{W})$.
\begin{enumerate}[label=(\roman*)]
\setlength{\itemsep}{0pt}
\item $(\ker(T^*T),\ker(TT^*))=(\ker(T),\ker(T^*))$.
\item $(\operatorname{range}(T^*T),\operatorname{range}(TT^*))=(\operatorname{range}(T^*),\operatorname{range}(T))$.
\item If $\mathcal{V}=\mathcal{U}\oplus\mathcal{Z}$, then $(\rho_{\mathcal{U},\mathcal{Z}})^*=\rho_{\mathcal{Z}^{\perp},\mathcal{U}^{\perp}}$.
\end{enumerate}
\begin{proof}
\begin{enumerate}[label=(\roman*)]
\setlength{\itemsep}{0pt}
\item Fix $\mathbf{v}\in\mathcal{V}$. Then $\mathbf{v}\in\ker(T^*T)\Longleftrightarrow0=\innerprod[\mathcal{V}]{\mathbf{v},T^*T\mathbf{v}}=\innerprod[\mathcal{W}]{T\mathbf{v},T\mathbf{v}}=\norm{T\mathbf{v}}^2\Longleftrightarrow\mathbf{v}\in\ker(T)$ by adjointness and definiteness.\\
The equality $\ker(TT^*)=\ker(T^*)$ follows by this result for $T^*$ and Proposition \ref{prop:basic-adjoint-identities}(iii).
\item The identities $(T^*T)^*=T^*T$ and $(TT^*)^*=TT^*$ hold by Proposition \ref{prop:basic-adjoint-identities}(ii),(iii).\\
The identity $\operatorname{range}(T^*T)=(\ker(T^*T))^{\perp}=(\ker(T))^{\perp}=\operatorname{range}(T^*)$ holds by (i) and Proposition \ref{prop:adjoint-kernel-range-orthogonality}(iii).\\
The identity $\operatorname{range}(TT^*)=(\ker(TT^*))^{\perp}=(\ker(T^*))^{\perp}=\operatorname{range}(T)$ holds by (i) and Proposition \ref{prop:adjoint-kernel-range-orthogonality}(iv).
\item Suppose $\mathcal{V}=\mathcal{U}\oplus\mathcal{Z}$ and set $P\defeq\rho_{\mathcal{U},\mathcal{Z}}$. Fix $\mathbf{c}\in\mathcal{Z}^{\perp}\cap\mathcal{U}^{\perp}$ and $\mathbf{x}\in\mathcal{V}$. Write $\mathbf{x}=\mathbf{u}+\mathbf{z}$ with $(\mathbf{u},\mathbf{z})\in\mathcal{U}\times\mathcal{Z}$.\\
Then $\innerprod{\mathbf{c},\mathbf{x}}=\innerprod{\mathbf{c},\mathbf{u}}+\innerprod{\mathbf{c},\mathbf{z}}=0$, so $\mathbf{c}=\mathbf{0}$ by Proposition \ref{prop:basic-properties-orthogonal-complement}(iii). Hence $\mathcal{Z}^{\perp}\cap\mathcal{U}^{\perp}=\{\mathbf{0}\}$.\\
The identity $\dim(\mathcal{Z}^{\perp}\oplus\mathcal{U}^{\perp})=\dim(\mathcal{V})$ follows by Corollaries \ref{cor:dimension-of-orthogonal-complement} and \ref{cor:dimension-direct-sum}, because $\dim(\mathcal{U})+\dim(\mathcal{Z})=\dim(\mathcal{V})$.\\
Hence $\mathcal{V}=\mathcal{Z}^{\perp}\oplus\mathcal{U}^{\perp}$ by Proposition \ref{prop:dimension-bounds-and-basis-criteria}(ii), and $Q\defeq\rho_{\mathcal{Z}^{\perp},\mathcal{U}^{\perp}}$ is defined.\\
Fix $\mathbf{x},\mathbf{y}\in\mathcal{V}$ and write $\mathbf{x}=\mathbf{u}+\mathbf{z}$ and $\mathbf{y}=\mathbf{a}+\mathbf{b}$ with $\mathbf{u}\in\mathcal{U}$, $\mathbf{z}\in\mathcal{Z}$, $\mathbf{a}\in\mathcal{Z}^{\perp}$, and $\mathbf{b}\in\mathcal{U}^{\perp}$.\\
By the projection formulas, $P\mathbf{x}=\mathbf{u}$ and $Q\mathbf{y}=\mathbf{a}$; by conjugate symmetry, $\innerprod{\mathbf{u},\mathbf{b}}=\innerprod{\mathbf{z},\mathbf{a}}=0$.\\
Consequently, $\innerprod{P\mathbf{x},\mathbf{y}}=\innerprod{\mathbf{u},\mathbf{a}+\mathbf{b}}=\innerprod{\mathbf{u},\mathbf{a}}=\innerprod{\mathbf{u}+\mathbf{z},\mathbf{a}}=\innerprod{\mathbf{x},Q\mathbf{y}}$. Thus $Q$ is the adjoint of $P$, and $P^*=Q$ by Lemma \ref{lem:adjoint-unique-if-exists}.
\end{enumerate}
\hlt{Structural Commentary}: Positivity fixes kernels, complementation fixes ranges, and decomposition fixes projection adjoints.\\
\hlt{Key Insight}: Adjoints preserve effective directions and orthogonally exchange projection complements.
\end{proof}
\end{proposition}

\begin{corollary} \label{cor:matrix-fundamental-subspaces-orthogonal-complements} \hlt{Matrix Fundamental Subspaces as Orthogonal Complements}\\
Let $A\in\operatorname{Mat}(m,n,\mathbb{F})$ act as $A:\mathbb{F}^n\to\mathbb{F}^m$ with the standard inner products. Then:
\begin{enumerate}[label=(\roman*)]
\setlength{\itemsep}{0pt}
\item $\ker(A)=\operatorname{col}(A^*)^{\perp}$;
\item $\operatorname{col}(A^*)=\ker(A)^{\perp}$;
\item $\ker(A^*)=\operatorname{col}(A)^{\perp}$;
\item $\operatorname{col}(A)=\ker(A^*)^{\perp}$.
\item $\mathbb{F}=\mathbb{R}\Longrightarrow\mathbb{R}^n=\ker(A)\odot\operatorname{row}(A)$.
\item $\mathbb{F}=\mathbb{R}\Longrightarrow\mathbb{R}^m=\ker(A^t)\odot\operatorname{col}(A)$.
\end{enumerate}
\begin{proof}
The induced map and its adjoint have ranges $\operatorname{col}(A)$ and $\operatorname{col}(A^*)$, respectively, by Theorem \ref{thm:matrix-of-adjoint-is-conjugate-transpose}.
\begin{enumerate}[label=(\roman*)]
\setlength{\itemsep}{0pt}
\item The identity $\ker(A)=\operatorname{col}(A^*)^{\perp}$ holds by Proposition \ref{prop:adjoint-kernel-range-orthogonality}(ii).
\item The identity $\operatorname{col}(A^*)=\ker(A)^{\perp}$ holds by Proposition \ref{prop:adjoint-kernel-range-orthogonality}(iii).
\item The identity $\ker(A^*)=\operatorname{col}(A)^{\perp}$ holds by Proposition \ref{prop:adjoint-kernel-range-orthogonality}(i).
\item The identity $\operatorname{col}(A)=\ker(A^*)^{\perp}$ holds by Proposition \ref{prop:adjoint-kernel-range-orthogonality}(iv).
\item If $\mathbb{F}=\mathbb{R}$, then $A^*=A^t$ and $\operatorname{col}(A^t)=\operatorname{row}(A)$. Hence $\mathbb{R}^n=\ker(A)\odot\operatorname{row}(A)$ by (i),(ii) and Theorem \ref{thm:orthogonal-direct-sum-decomposition}.
\item If $\mathbb{F}=\mathbb{R}$, then $A^*=A^t$. Hence $\mathbb{R}^m=\ker(A^t)\odot\operatorname{col}(A)$ by (iii),(iv) and Theorem \ref{thm:orthogonal-direct-sum-decomposition}.
\end{enumerate}
\hlt{Structural Commentary}: Standard coordinates convert operator formulas into matrix fundamental-subspace identities.\\
\hlt{Key Insight}: Row and column spaces are orthogonal complements of their nullspaces.
\end{proof}
\end{corollary}
